% !TeX TS-program = xelatex
% !Mode TeXUTF-8
%# -- codingutf-8 --
% !TeX encoding = UTF-8
%%%%%%%%%%%%%%%%%%%%%%%%%%%%%%%%%%%%%%%%%
% 研究意义、研究内容、技术路线、研究方法和进度在下面填写
%--------论文评定成绩---------------------
\newcommand{\yiyi}{
%------------研究意义----------------------
\setlength{\parindent}{2em}
%\vspace{2mm}		
	矩阵正定是高等代数的重要内容，其判定方法是矩阵理论研究的重要分支。正定矩阵具有独特的代数性质和几何意义，在优化理论、数值分析、概率论等多个数学分支中扮演着基础支撑角色。矩阵正定的判定方法繁多，但各类方法散落在高等代数教材当中，缺乏系统的归纳整理，导致在实际应用中难以快速筛选合适的方法计算。

	本课题基于这种现象， 以归纳整理矩阵正定的判定方法为出发点，系统整合现有判定体系，厘清不同方法的内在关联与适用边界，为矩阵理论的体系化发展提供支撑， 同时为复杂矩阵的正定判定提供更清晰的理论指引
%-------------------------------------------
}
\newcommand{\neirong}{
	\setlength{\parindent}{2em}
%----------------------研究内容---------------------
本课题以矩阵正定的判定条件及应用为核心，涵盖多方面内容，深入分析矩阵正定的性质及判定条件， 旨在厘清不同判定条件的适用范围，为矩阵理论的体系化发展提供支撑， 同时为复杂矩阵的正定判定提供更清晰的理论指引。

在引言中，结合高等代数学科发展史及正定矩阵在高等代数领域的需求， 阐明矩阵正定判定方法及应用的重要性。

理论基础部分将系统梳理矩阵正定的基本概念与相关性质，为后续判定方法的研究奠定基础。

案例分析部分将选取高等代数教材的经典例题，运用上述判定方法进行实证分析。通过案例研究验证各判定方法的有效性与实用性，总结不同场景下判定方法的选择技巧。

结论部分归纳整理矩阵正定的判定条件及应用，总结课题研究过程的收获与不足，此外还将列出参考文献并撰写致谢。

%----------------------------------------------------
}
\newcommand{\luxian}{
%------------技术路线、研究方法和进度------------------
\setlength{\parindent}{2em}
3.1 技术路线和研究方法

本研究首先系统梳理矩阵正定相关的理论知识，其次，全面收集国内外关于矩阵正定判定方法及应用的研究文献，进行分类整理与深入分析。随后，选取典型例题，运用相应判定方法进行实证分析，验证方法的有效性并总结应用技巧。最后， 总结研究成果，分析不足并展望未来研究方向，完成论文撰写与修改。研究方法采用文献研究法和案例分析法。文献研究法用于查阅资料提供理论支持，案例分析法通过剖析案例总结方法。\\
3.2 研究进度\\
第七学期\\
2025/11/10-2026/1/10 ：确定具体题目、撰写提纲和开题报告\\
2026/1/10前：学生网上提交开题报告\\
第八学期\\
2026/3/15前：继续搜集文献资料、素材，撰写初稿\\
2026/3/23前： 网上提交初稿\\
2026/3/24-4/13 ：在导师指导下修改初稿\\
2026/4/28前： 网上提交定稿\\

%----------------------------------------------------
}